\documentclass{article}
\usepackage{ifthen,array,booktabs}
\usepackage{ctex}
\usepackage{tikz,xcolor,colortbl}
\usepackage[left=1in,top=1.1in,right=1in,bottom=1.1in]{geometry}
\usepackage{hyperref}
\usepackage{fontspec}

\definecolor{Cyan}{RGB}{204,231,224}
\definecolor{Pink}{RGB}{252,228,208}
\definecolor{dpink}{RGB}{238,202,183}
\definecolor{Violet}{RGB}{232,232,237}
\definecolor{Orange}{RGB}{255,228,180}
\definecolor{Blue}{RGB}{173,213,237}


\IfFileExists{csvsimple-doc.cfg}{\input{csvsimple-doc.cfg}}{}% ignore this line
\providecommand\csvpkgprefix{}                               % ignore this line

\usepackage{\csvpkgprefix csvsimple-l3}%  \usepackage{csvsimple-l3}

\begin{document}
	
	\title{\bfseries\sffamily CSV.example for Automatic Table Generation}
	
	\author{\href{http://www.latexstudio.net}{\sffamily csvsimple-l3 宏包自动生成CSV数表表格的几个案例}}
	
	\date{September 16, 2021}
	
	\maketitle
	
	%----------------------------------------------------------
	\section{CSV 数表文本自动生成（用于以下信息读取测试）}
	我们先建立一个csv格式的数据总表，并将其在\LaTeX 文档中被读取出来。
	
	\smallskip
	{\small
		\csvautotabular{csvsimple-example.csv}
	}
	\smallskip
	
	以下自动生成的表内容，均源于宏包csvsimple-l3对上面总表信息内容的读取。
	
	%----------------------------------------------------------
	\section{读取数据生成的第一个表格}
	\csvreader[
	tabular=|l|l|l|l|,
	table head=\hline\multicolumn{4}{|c|}{\cellcolor{dpink}\bfseries 高一(2)班学生语数外成绩表}\\\hline
	\cellcolor{Blue}\bfseries 姓名 &\cellcolor{Blue} \bfseries 语文&\cellcolor{Blue} \bfseries 数学&\cellcolor{Blue} \bfseries 英语\\\hline\hline,
	late after line=\\\hline]
	{csvsimple-example.csv}{Name=\Name,Chinese=\Chinese,Math=\Math,
		English=\English}{
		\cellcolor{Violet}\Name& \cellcolor{Orange}\Chinese & \cellcolor{Pink}\Math & \cellcolor{Cyan}\English
	}\hspace{2cm}
	\csvreader[tabular=|l|l|l|l|,
	table head= \hline\multicolumn{4}{|c|}{\cellcolor{red!70!yellow!80!black}\color{white}\bfseries 高一(2)班学生理化生成绩表}\\\hline
	\cellcolor{violet!70!cyan!80!white}\color{white}\bfseries 姓名 & \cellcolor{violet!70!cyan!80!white}\color{white}\bfseries 物理& \cellcolor{violet!70!cyan!80!white}\color{white}\bfseries 化学& \cellcolor{violet!70!cyan!80!white}\color{white}\bfseries 生物\\\hline\hline,
	late after line=\\\hline]
	{csvsimple-example.csv}{Name=\Name,Physics=\Physics,Chemical=\Chmical,
		Biology=\Biology}{
		\cellcolor{Blue!70!green!80!white}\Name& \cellcolor{pink!70!red!80!white}\Physics & \Chmical & \Biology
	}
	
	
	%----------------------------------------------------------
	\section{读取数据，自动生成格式一样的文档}
	\csvnames{my names}{Name=\Name,Chinese=\Chinese,Math=\Math,
		English=\English,Total=\Total,Overall=\Overall}
	
	\csvreader[my names, late after line=, late after last line=]%
	{csvsimple-example.csv}{Color=\Color}{
		\begin{minipage}{0.8\linewidth}
			\begin{tikzpicture}[overlay,remember picture]
				\draw[line width=1pt,draw=orange,fill=\Color,dashed,rounded corners=6pt](-1,-5) rectangle (15,0);
			\end{tikzpicture}
			\begin{center}
				{\LARGE \bfseries 奖\quad~状}
			\end{center}
			
			\noindent
			\underline{\Name}同学：
			
			在2020--2021学年度第一学期中期末考试中，语文: \underline{~\Chinese~}、数学: \underline{~\Math~}、英语\underline{~\English~}、总分:\underline{~\Total~}，总体表现 \underline{\Overall}，特发此奖以资鼓励！
			
			\vspace{15pt}
			\hfill XXXX 中学 \hspace{2cm}
			
			\hfill 2021年9月 16日\hspace{2cm}
		\end{minipage}
		
		\vspace{1cm}
	}
	
\end{document}

